% Part: first-order-logic
% Chapter: natural-deduction
% Section: provability-quantifiers

% verification of properties of provability needed for maximally
% consistent sets in the completeness chapter.

\documentclass[../../../include/open-logic-section]{subfiles}

\begin{document}

\olfileid{fol}{ntd}{qpr}

\olsection{\usetoken{S}{derivability} અને પરિમાણકો}

\begin{explain}
  પૂર્ણતા-પ્રમેય માટે એ પણ જરૂરી છે કે સ્વાભાવિક નિગમનના નિયમો આ
  વિભાગમાં સ્થાપેલી~$\Proves$ વિશેની હકીકતો આપે.
\end{explain}

\begin{thm}
\ollabel{thm:strong-generalization} જો $c$ એવું અચલ હોય જે $\Gamma$ અથવા
$!A(x)$માં આવતું ન હોય અને $\Gamma \Proves !A(c)$ હોય, તો
$\Gamma \Proves \lforall[x][!A(x)]$.
\end{thm}

\begin{proof}
$\Gamma$માંથી $!A(c)$ની !!a{derivation}ને $\delta$ કહો.
\Intro{\lforall} અનુમાન ઉમેરવાથી $\lforall[x][!A(x)]$ની
!!a{derivation} મળે છે. $c$ એ $\Gamma$ કે $!A(x)$માં આવતું નથી, તેથી
આઇગનચલ-શરત સંતોષાય છે.
\end{proof}

\begin{prop}
\ollabel{prop:provability-quantifiers}
\begin{tagenumerate}{prvEx,prvAll}
\tagitem{prvEx}{$!A(t) \Proves \lexists[x][!A(x)]$.}{}

\tagitem{prvAll}{$\lforall[x][!A(x)] \Proves !A(t)$.}{}
\end{tagenumerate}
\end{prop}

\begin{proof}
\begin{tagenumerate}{prvEx,prvAll}
\tagitem{prvEx}{નીચે આપેલું~$!A(t)$માંથી~$\lexists[x][!A(x)]$ની
  !!a{derivation} છે:
\begin{prooftree}
\AxiomC{$!A(t)$}
\RightLabel{\Intro{\lexists}}
\UnaryInfC{$\lexists[x][!A(x)]$}
\end{prooftree}}{}

\tagitem{prvAll}{નીચે આપેલું~$\lforall[x][!A(x)]$માંથી~$!A(t)$ની
  !!a{derivation} છે:
\begin{prooftree}
\AxiomC{$\lforall[x][!A(x)]$}
\RightLabel{\Elim{\lforall}}
\UnaryInfC{$!A(t)$}
\end{prooftree}}{}
\end{tagenumerate}
\end{proof}

\end{document}
